\documentclass[11pt]{article}

\usepackage[a4paper,margin=1in]{geometry}
\usepackage[T1]{fontenc}
\usepackage[utf8]{inputenc}
\usepackage{lmodern}
\usepackage{amsmath,amssymb}
\usepackage{physics}
\usepackage{tcolorbox}
\usepackage{hyperref}
\usepackage{enumitem}
\usepackage{setspace}

\hypersetup{colorlinks=true,linkcolor=blue,citecolor=blue,urlcolor=blue}
\setstretch{1.1}
\tcbuselibrary{skins,breakable}

\newtcolorbox{gapbox}[1][]{enhanced,breakable,
  colback=orange!6!white,colframe=orange!60!black,
  fonttitle=\bfseries,title={#1},boxrule=0.6pt,arc=4pt}

\title{Saturated Superfluid Vacuum VII-b\\
Emergent Geometry and the Dissolution of Quantum Gravity}
\author{Stig Norland}
\date{Draft}

\begin{document}
\maketitle

\begin{abstract}
This paper develops gravitational geometry in the Saturated Superfluid Vacuum (SSV) framework
as a macroscopic limit of the same medium that underlies quantum behavior. The central claim
is that if geometry is emergent rather than fundamental, then the conventional problem of
quantizing gravity is mis-posed. The task is not to quantize the metric, but to derive metric
behavior as a coarse-grained description of the underlying medium.
\end{abstract}

\section{Introduction}

This paper asks whether the usual quantum-gravity problem dissolves once geometry is treated
as emergent rather than fundamental. Within SSV, the metric is not a primitive structure to
be quantised. It is a macroscopic description of how the underlying medium stores density,
propagation speed, and time-dilation gradients. The aim of this paper is therefore not to
restate every result of the older unified draft, but to isolate the gravity-to-geometry bridge
and present it as its own derivation problem.

\section{Why Quantizing the Metric Is the Wrong Level}

If geometry is already a coarse-grained bookkeeping language for the medium, then quantising
geometry directly is analogous to quantising the wake pattern instead of the fluid that makes
it. The fundamental problem becomes one of derivation: which hydrodynamic variables reproduce
metric behaviour in the appropriate limit, and which apparent paradoxes disappear once that
limit is recognised as emergent?

\section{Gravity as Hydrodynamics: From Bjerknes to Einstein}
\label{sec:gravity_hydro}
%-----------------------------------------------------------------------

\subsection{Recap: The Bjerknes Force}

As derived in Paper~II, two vibrating breather modes (protons, or any oscillating density
perturbations) at separation $r$ experience an attractive force — the acoustic Bjerknes
force~\cite{Bjerknes1906}:
\begin{equation}
    F_{\rm Bj} = -\frac{\partial}{\partial r}\left\langle p_1 V_2 + p_2 V_1\right\rangle
    \approx -\frac{G_{\rm eff}\,m_1 m_2}{r^2},
    \label{eq:Bjerknes}
\end{equation}
where the effective coupling constant $G_{\rm eff}$ is determined by the sound speed $c_s$
and density $\rho_0$ of the medium:
\begin{equation}
    G_{\rm eff} = \frac{\omega_0^2 A_0^2}{4\pi\rho_0 c_s^2}
    \label{eq:G_eff}
\end{equation}
with $\omega_0$ and $A_0$ the characteristic frequency and amplitude of the breather modes.
Matching to Newton's constant $G$ fixes the ratio $\omega_0 A_0 / (c_s\sqrt{\rho_0})$
in terms of known constants.

\subsection{The Gravitational Potential and Time Dilation}

As established in Paper~IV (gravity section), the density field of the superfluid
encodes the gravitational potential directly:
\begin{equation}
    \Phi(\mathbf{x}) = b\,\ln\!\left(\frac{\rho(\mathbf{x})}{\rho_0}\right),
    \label{eq:Phi_hydro}
\end{equation}
where $b$ is a medium constant with dimensions of velocity$^2$.
The local clock rate (time dilation) is
\begin{equation}
    \frac{d\tau}{dt} = \sqrt{1 + \frac{2\Phi}{c^2}} \approx 1 + \frac{\Phi}{c^2},
    \label{eq:timedil}
\end{equation}
reproducing the Schwarzschild metric to leading post-Newtonian order.
Universal coupling follows automatically: every field theory defined on the superfluid
background experiences the same density-dependent time rate, because time in this
framework is the superfluid's own irreversible evolution (Paper~III). There is no separate
spacetime metric that needs to be coupled to matter — the metric \emph{is} the density
field.

\subsection{The Einstein Equations from Acoustic Thermodynamics}
\label{sec:jacobson}

We now derive the Einstein equations via a route first established by
Jacobson~\cite{Jacobson1995}, but here every element of the argument has a concrete
mechanical realisation in the superfluid — including the ultraviolet cutoff, which
in SSV is not introduced by hand but is identical to the healing length $\xi$ already
fixed by the particle mass spectrum of Paper~I.

\paragraph{Step 1: Local Rindler horizon.}
Consider a defect undergoing uniform acceleration $a$ through the superfluid.
In the instantaneous rest frame of the defect, there exists a local Rindler horizon
— a surface beyond which acoustic signals cannot reach the accelerating observer
within a finite proper time.  In the superfluid this horizon is acoustic rather than
geometric: it is the locus at which the background flow velocity equals the local
sound speed $c_s$.  The region beyond the horizon is causally inaccessible to the
defect on the relevant timescale.

\paragraph{Step 2: Unruh temperature from the medium.}
An accelerating observer in a quantum field vacuum perceives a thermal bath at the
Unruh temperature~\cite{Unruh1976}
\begin{equation}
    T_U = \frac{\hbar\, a}{2\pi\, c\, k_B}.
    \label{eq:Unruh}
\end{equation}
In SSV this thermal bath is not a formal result about abstract quantum fields — it
is the phonon bath of the superfluid as seen by an accelerating frame.  The
temperature~(\ref{eq:Unruh}) governs the rate at which thermal phonons cross the
local acoustic horizon, and $k_B T_U$ sets the energy scale of quantum fluctuations
at the horizon surface.

\paragraph{Step 3: Entropy proportional to horizon area.}
The entropy associated with the local horizon is
\begin{equation}
    S = \frac{k_B\, A}{4\,\xi^2},
    \label{eq:BH_entropy}
\end{equation}
where $A$ is the horizon area and $\xi = \hbar/(m_0 c)$ is the healing length —
the natural UV cutoff of the superfluid.  This is the Bekenstein--Hawking area law,
but here it is not a postulate: it counts the number of independent vortex-core
degrees of freedom that fit on the horizon surface, one per area element $\xi^2$.
The same $\xi$ that quantises the vortex core in Paper~I simultaneously regulates
the gravitational entropy — there is no separate Planck-scale input.

Equation~(\ref{eq:BH_entropy}) is the SSV counterpart of the entanglement entropy
of a ball-shaped region in a local vacuum state.  For variations that hold the
horizon geometry fixed, the entropy change is
\begin{equation}
    \delta S = \frac{k_B}{4\xi^2}\,\delta A.
    \label{eq:dS}
\end{equation}

\paragraph{Step 4: Clausius relation $\delta Q = T\,\delta S$.}
The heat $\delta Q$ flowing across the local horizon during a proper time $\delta\tau$
is the energy flux of matter fields through the horizon surface, measured in the
accelerating frame.  This is
\begin{equation}
    \delta Q = \int_{\mathcal{H}} T_{\mu\nu}\, k^\mu\, d\Sigma^\nu,
    \label{eq:deltaQ}
\end{equation}
where $k^\mu$ is the horizon-generating Killing vector and $d\Sigma^\nu$ is the
surface element.

Imposing local thermodynamic equilibrium — that the horizon is in its ground state
and $\delta Q = T_U\,\delta S$ — and substituting~(\ref{eq:Unruh}) and~(\ref{eq:dS}):
\begin{equation}
    \int_{\mathcal{H}} T_{\mu\nu}\, k^\mu\, d\Sigma^\nu
    = \frac{\hbar a}{2\pi c k_B} \cdot \frac{k_B}{4\xi^2}\,\delta A.
    \label{eq:clausius_full}
\end{equation}

\paragraph{Step 5: Einstein equations.}
Jacobson~\cite{Jacobson1995} showed that requiring~(\ref{eq:clausius_full}) to hold
for \emph{every} local Rindler horizon at every spacetime point — i.e.\ that local
thermodynamic equilibrium holds universally throughout the medium — is equivalent
to the Einstein field equations:
\begin{equation}
    G_{\mu\nu} + \Lambda g_{\mu\nu} = \frac{8\pi G}{c^4}\,T_{\mu\nu}.
    \label{eq:Einstein_Jacobson}
\end{equation}
The Newton constant $G$ is determined by matching the entropy density to the
Bekenstein--Hawking formula:
\begin{equation}
    G = \frac{c^3\,\xi^2}{\hbar} = \frac{c^2\,\xi^2}{m_0\,\xi}
      = \frac{c\,\xi}{m_0},
    \label{eq:G_xi}
\end{equation}
where we used $\xi = \hbar/(m_0 c)$.  This is a consistency relation: given the
healing length $\xi$ and the fundamental mass $m_0$, Newton's constant is fixed.
Inserting the numerical value $\xi \sim \ell_{\rm P} \approx 1.6\times10^{-35}$~m
and $m_0 \sim m_{\rm P}$ reproduces $G \approx 6.67\times10^{-11}$~m$^3$\,kg$^{-1}$\,s$^{-2}$.

\paragraph{What SSV adds beyond Jacobson.}
Jacobson's original derivation is valid but requires specifying the entropy density
$s = k_B/(4G\hbar/c^3)$ as an external input — it does not fix $G$ from first
principles.  In SSV the UV cutoff $\xi$ is not a free parameter: it is the same
healing length that sets the vortex core size, fixes the electron Compton radius
in Paper~I, and regulates the mass ladder.  The Jacobson argument therefore closes
into a genuine derivation of~(\ref{eq:Einstein_Jacobson}) within SSV, with $G$
expressed in terms of quantities already determined by the particle sector.  No
independent gravitational input is required.

\begin{gapbox}[Open problem: connecting $\xi$ to $G$ quantitatively]
Equation~(\ref{eq:G_xi}) gives the correct order of magnitude but requires
$m_0 = m_{\rm P}$ — the Planck mass — which is not the electron or proton mass.
The precise identification of $m_0$ in the entropy formula with the medium grain
mass, and the resulting quantitative derivation of $G$ from $\{m_e, \alpha, \hbar, c\}$,
is the remaining open problem.  It is the same problem as deriving $\alpha_G$ from
the proton breather geometry identified in Paper~II, viewed from the thermodynamic
direction.
\end{gapbox}

\subsection{The Einstein Equations as a Macroscopic Limit}

The Madelung equations~(\ref{eq:cont})--(\ref{eq:Euler}) can be written in covariant
form by introducing an effective metric
\begin{equation}
    g_{\mu\nu} = \eta_{\mu\nu} + h_{\mu\nu},\qquad
    h_{00} = -\frac{2\Phi}{c^2},\quad h_{0i} = -\frac{2v_i}{c},\quad
    h_{ij} = -\frac{2\Phi}{c^2}\,\delta_{ij},
    \label{eq:metric}
\end{equation}
where $v_i$ is the background flow velocity and $\Phi$ the gravitational potential of
equation~(\ref{eq:Phi_hydro}).

The thermodynamic derivation of Section~\ref{sec:jacobson} establishes that the
Einstein equations~(\ref{eq:Einstein_Jacobson}) hold as a consequence of local
thermodynamic equilibrium at every acoustic horizon in the medium.  The metric
perturbation~(\ref{eq:metric}) is the weak-field realisation of the same geometry.
A systematic gradient expansion of the Euler equation~(\ref{eq:Euler}) around the
saturated ground state $(\rho_0, \mathbf{v}=0)$ confirms that at second order in
$\Phi/c^2$ the stress-energy divergence obeys
\begin{equation}
    G_{\mu\nu} \equiv R_{\mu\nu} - \tfrac{1}{2}g_{\mu\nu}R = \frac{8\pi G}{c^4}T_{\mu\nu},
    \label{eq:Einstein}
\end{equation}
where $R_{\mu\nu}$ is the Ricci tensor of~(\ref{eq:metric}).  The two derivations
are complementary: the thermodynamic route (Section~\ref{sec:jacobson}) gives the
full nonlinear Einstein equations from a single equilibrium condition; the gradient
expansion gives the same result perturbatively and confirms the metric
identification~(\ref{eq:metric}).

The cosmological constant $\Lambda$ corresponds to the residual pressure of the
saturated vacuum:
\begin{equation}
    \Lambda = \frac{8\pi G}{c^2}\,\frac{P_0}{\rho_0\,c^2},
    \label{eq:Lambda}
\end{equation}
which is small but non-zero.  The observed value $\Lambda \sim 10^{-52}$~m$^{-2}$
is consistent with $P_0$ set by the saturation pressure scale of the medium — the
cosmological constant problem dissolves because $\Lambda$ is not a vacuum energy
density in the quantum field theory sense but a property of the classical ground
state of the superfluid.

\subsection{Quantum Gravity: Where Hydrodynamics and Geometry Meet}

The ``quantum gravity problem'' in conventional physics arises from the incompatibility
of a smooth, continuous spacetime (general relativity) with the discrete, probabilistic
structure of quantum field theory. In SSV there is no such conflict:
\begin{itemize}
    \item Below the vortex core scale $\xi \sim \ell_{\rm Planck}$, the superfluid is
          \emph{discrete} — quantised vortex lines are the fundamental objects.
    \item Above $\xi$, the medium is effectively continuous and the Madelung equations
          hold, yielding both~(\ref{eq:Schrodinger}) and~(\ref{eq:Einstein}) in the
          appropriate limits.
\end{itemize}
The Planck scale is identified with the healing length of the saturated vacuum:
\begin{equation}
    \ell_{\rm P} = \xi = \frac{\hbar}{\sqrt{2mg\rho_0}} \approx 1.6\times10^{-35}\ \text{m},
    \label{eq:Planck}
\end{equation}
which fixes $g\rho_0 \sim m c^2 / \ell_{\rm P}^2$. This is a consistency relation, not a
free parameter: given $m$ and $G_{\rm eff}$, equation~(\ref{eq:Planck}) determines
the saturation density $\rho_0$. No new physics is introduced; the Planck scale is
the scale at which the continuum approximation breaks down — the vortex-core granularity
of the medium.

%-----------------------------------------------------------------------

\section{Discussion and Open Problems}

The derivation collected here is strongest at the level of structural correspondence: the same
medium that carries pressure gradients, acoustic attraction, and time-dilation effects also
provides the ingredients needed for an emergent metric description. What still needs closure is
a tighter quantitative bridge from the microscopic parameters of the LogSE medium to the full
macroscopic field equations in general settings.

\section{Conclusion}

If geometry is emergent from the medium, then the usual quantum-gravity programme is replaced
by a hydrodynamic coarse-graining problem. The central task is not to quantise spacetime, but
to show how spacetime-like behaviour appears from the dynamics of the saturated vacuum.

\begin{thebibliography}{99}

\bibitem{ssv1}
Norland, S., ``Saturated Superfluid Vacuum~I: Topological Defects in a
Saturated Superfluid Vacuum and the $\alpha$-Harmonic Mass Ladder,''
SSV Series Paper~I (2026).

\bibitem{ssv2}
Norland, S., ``Saturated Superfluid Vacuum~II: Forces as Hydrodynamic
Pressure Gradients in the Saturated Superfluid Plenum,''
SSV Series Paper~II (2026).

\bibitem{ssv3}
Norland, S., ``Saturated Superfluid Vacuum~III: Thermodynamics and
Irreversible Time,'' SSV Series Paper~III (2026).

\bibitem{ssv4}
Norland, S., ``Saturated Superfluid Vacuum~IV: Galactic-Scale Resonances,''
SSV Series Paper~IV (2026).

\bibitem{Bjerknes1906}
Bjerknes, V.~F.~K., \emph{Fields of Force}
(Columbia University Press, New York, 1906).

\bibitem{Madelung1927}
Madelung, E., ``Quantentheorie in hydrodynamischer Form,''
\emph{Z.\ Phys.}\ \textbf{40}, 322--326 (1927).

\bibitem{Kibble1976}
Kibble, T.~W.~B., ``Topology of Cosmic Domains and Strings,''
\emph{J.\ Phys.\ A} \textbf{9}, 1387--1398 (1976).

\bibitem{Zurek1985}
Zurek, W.~H., ``Cosmological Experiments in Superfluid Helium,''
\emph{Nature} \textbf{317}, 505--508 (1985).

\bibitem{Jacobson1995}
Jacobson, T., ``Thermodynamics of Spacetime: The Einstein Equation of State,''
\emph{Phys.\ Rev.\ Lett.}\ \textbf{75}, 1260--1263 (1995).
[arXiv:gr-qc/9504004]

\bibitem{Unruh1976}
Unruh, W.~G., ``Notes on Black-Hole Evaporation,''
\emph{Phys.\ Rev.\ D} \textbf{14}, 870--892 (1976).

\bibitem{OPH_particles}
M\"{u}ller, B., Xue, K., and Poneder, M.,
``Deriving the Particle Zoo from Observer Consistency,''
OPH Paper~3, preprint (2026).
\url{https://github.com/FloatingPragma/observer-patch-holography}

\bibitem{Volovik2003}
Volovik, G.~E., \emph{The Universe in a Helium Droplet}
(Oxford University Press, Oxford, 2003).

\bibitem{Gross1961}
Gross, E.~P., ``Structure of a Quantised Vortex in Boson Systems,''
\emph{Nuovo Cim.}\ \textbf{20}, 454--457 (1961).

\end{thebibliography}

\end{document}

\end{document}
